% v005 (2026-11-17): introduction for the "empirical study" narrative. First paragraph kept from v003 (citations
% unchanged); positioning follows chapters/related_work/positioning.tex; numbers from the Results section.
\section{Introduction}
\label{sec:introduction}

Evaluating a candidate network against the true objective -- by training it, through weight sharing, or by
querying a benchmark -- dominates the cost of neural architecture search~\citep{white2023insights,liu2023survey}.
Architecture and training hyperparameters are commonly searched in sequence, yet Zela et al.\ show that they
interact, so the best architecture under one hyperparameter setting need not be the best under
another~\citep{zela2018towards}. Searching both jointly avoids that risk at the price of a larger space, in which
some variables matter only together.

A family of evolutionary algorithms is built around exactly that structure. GOMEA and the Parameter-less Population
Pyramid (P3) learn a linkage tree of dependent variables from the population and keep those groups together during
crossover~\citep{thierens2011optimal,goldman2014parameterless}; on combinatorial benchmarks their advantage over
structure-blind search grows with problem size~\citep{goldman2015fast,qiao2024gomearuntime}. Members of the family
have been applied to architecture search~\citep{denottelander2021localsearch,chebykin2022encas,dahi2026tightening},
and surrogate-assisted variants exist~\citep{dushatskiy2021novelsurrogateassistedevolutionaryalgorithm,przewozniczek2026lympus}.
What is missing is a measurement: how P3 compares, on the joint multi-objective task, with the optimisers that are
actually used for it, at what budget, and at what cost.

\paragraph{This study.} We compare P3 in five forms---canonical P3; P3Net, which replaces the evaluation of mixing steps
by a pairwise surrogate of fitness differences; two variants of P3Net that isolate how the surrogate is used and what
the linkage tree adds to it; and P3 with the eLyMPuS surrogate~\citep{przewozniczek2026lympus}---with 22 optimisers
from NAS, evolutionary computation, Bayesian optimisation and multi-fidelity optimisation. The benchmarks are
JAHS-Bench-201~\citep{bansal2022jahsbench} on three tasks and NAS-HPO-Bench-II~\citep{hirose2021nashpobenchii}, with
four FCNet tasks~\citep{klein2019tabular} that played no part in designing P3Net and NAS-Bench-201 with an energy
objective~\citep{dong2020nasbench201,li2021hwnasbench}. Methods, metrics, tests and seeds were fixed in a
pre-registration before the confirmatory runs, and every sentence about a method's behaviour is tied to a test of the
implementation.

\paragraph{Findings.} P3's standing is set by the budget. With 50 or 100 evaluations it is worse than most arms on
JAHS-Bench-201, random search included, because it spends them on hill-climbing single solutions before any linkage
is used. With 200 or 350 it is worse than no arm on two of the three JAHS-Bench-201 tasks and returns fronts that are
strongest at the accurate end, while population-based methods stay ahead on NAS-HPO-Bench-II. The surrogate changes
nothing that our tests can detect, also at 1\,000 evaluations where it screens most mixing steps, and it raises the
optimiser's own time by an order of magnitude. On the four FCNet tasks P3Net wins more comparisons than it loses.

\paragraph{Contributions.}
\begin{itemize}
    \item A pre-registered comparison of a linkage-learning optimiser with 22 optimisers for joint multi-objective
    NAS and hyperparameter optimisation, on nine search spaces, four budgets and 30 seeds per cell, with
    Holm--Bonferroni-corrected tests, effect sizes, equivalence tests and the optimiser's own time and memory.
    \item A controlled account of what a pairwise surrogate adds to P3: ablations that isolate replacement versus
    gating and the role of the linkage tree, a check that the surrogate-free P3Net is exactly canonical P3, and an
    exploratory budget at which the surrogate acts on most mixing steps.
    \item Evidence on the benchmarks themselves: retraining the selected architectures from scratch, and secondary
    analyses (anytime quality, cost to target, regret, generalisation gaps, sensitivity of the metric).
    \item Code, run records, analysis scripts and the pre-registration, with every table regenerated from the
    records.
\end{itemize}
